\section{The coefficient chain at an entirely thick inert place}

The companion \emph{RH Lean Roadmap}, revision
\texttt{c9112b83e6fe}, gives the coefficient-chain construction
in B8.42bu--B8.42bu1, with its longer proofs in the same volume.
We use its notation locally and give the integral jet and content
arguments on which the comparison of exponents depends.

% BEGIN T83A1
\paragraph{T83A1. The all-thick data and the two Frobenius truncations.}
Take a guarded inert C or D index, a good point, a place whose
word is entirely thick, any good coefficient extension and any tested
crystalline surjection $F:T_1\to T_2$. By I82W1--I82W5 and C82S1
the completion is unramified of degree $n=2$ and $p\ge11$.
Set $q=p-1$, $r=\operatorname{rank}T_2$. The large jump intervals
are $[-\beta_i-S,-\beta_i]$ at the negative letter and
$[\beta_i,\beta_i+S]$ at its neighbour, with the original parity.
The remaining jumps lie in $[-\delta,\delta]$ and their
multiplicity on the target is at most $m$ at each letter.
Strictness gives the same enclosure on the target. Write
\[
 h=\max_i\beta_i+S,\quad H=h+\delta,\quad
 \lambda=1+\log_p(2H),\quad
 a=\left\lceil H(2p-1)/q^2+3+\lambda\right\rceil .
\]
The thickness inequalities give $h/2\le\beta_i$ and
$2a+\delta\le\beta_i+1$; in fact $2a+\delta\le\beta_i$
when $2a+\delta$ is required in an interval separation, since
$2\lceil A\rceil<2A+2$ and $2A+\delta+1\le\beta_i$.
Use the invariant Wach basis over $A_0=O[[v]]$, $Q=v+p$,
$Q_1=\varphi Q$ and $u=1+p$. In this basis
$\varphi(v)\sim vQ^q$ and $Q_1=p+vQ^qU$ for a unit $U$.
The least power annihilator and the descending divided difference
of R82W1--R82W2 apply to $F$ and to each crystalline extension
surjecting onto the trivial lattice. They imply that $Q^a$ kills
each Wach cokernel. Indeed $v^j\sim X^{qj}$ in the original
variable, so the term $p^{a-j}v^j$ in $(p+v)^a$ suffices when
$qj\ge H$. Otherwise R82W2 asks for
$a-j\ge(H-qj)p/q^2+\lceil1+\log_p(H+1)\rceil$.
This follows already from $a\ge Hp/q^2+\lambda+1$.
Averaging over the tame group gives invariant lifts, including the
lifts of a fixed scalar $p^N$ at all letters from R82W2 with $j=0$.

For integers $M,C\ge0$ with $M\ge1$, a normalized constant-period
frame and its inverse can each be represented integrally after
$\varphi$ modulo $v^M Q^C$ at cost
$s=\lceil(M+C)/q+1\rceil$. After $\varphi^2$ modulo
$vQ^C$ the cost is $s_2=\lceil C/q^2+1\rceil$. Here the statement
means that $p^s$ times the substituted matrix, or $p^{s_2}$
respectively, has integral polynomial-jet coordinates in the
translated basis. For the first assertion the degree $k$ input
has denominator at most $p^{w(k)}$, $w(k)=k+v_p(k!)\le pk/q$,
by J82F1. Its image contains $v^kQ^{qk}$. If $k\ge M$ and
$qk\ge C$, it vanishes in the quotient. If only $k\ge M$,
remove $v^M$ and expand $(Q-p)^{k-M}$; if only $qk\ge C$,
remove $Q^C$ and expand $(v+p)^{qk-C}$. In either case every
term surviving modulo $v^M Q^C$ gains at least
$\max(0,pk-M-C)$ factors of $p$. If neither inequality holds,
$pk<M+C$. The resulting valuation covers $w(k)-s$ in every case.

For the second assertion, a term $k\ge1$ contains
$v^kQ^{qk}Q_1^{qk}$. In expanding the last factor, any term
surviving modulo $vQ^C$ uses $j$ copies of $vQ^qU$ with
$q(k+j)<C$, and hence gains $qk-j>pk-C/q$ factors of $p$.
If $pk\le C/q$, $w(k)\le C/q^2$; above this threshold the just
computed gain grows faster than $w(k)$. The constant term is
unchanged. Both calculations need only finitely many $k$.
The two matrices multiply to the corresponding scalar square in
the quotient, by the same polynomial division and $p$-saturation
as J82F1. They identify the Frobenius basis with its constant
crystalline periods at every primitive branch to order $C$,
and give the prescribed constant at $v=0$. Indeed after the
indicated Frobenius steps these expansions are taken at $v=0$
before substitution, exactly as in J82F2.
% END T83A1

% BEGIN T83A2
\paragraph{T83A2. The constant coefficient gains the full negative-letter saving.}
Let $c$ be negative, $b=\beta_c$, $y=\beta_{c-1}+S$,
$X_0=2b+S+\delta$ and $s_2=\lceil X_0/q^2+1\rceil$.
Write $N_i$ for the invariant Wach component of the target and
$R_i^{(j)}=A_0\varphi^jN_i$. Take any vector $w$ at $c$ with
$(\gamma-1)w\in vN_c$, $Q^\delta w\in R_{c-1}^{(1)}$
and finite Wach poles. The one-step lower height implies
$Q^{\delta+b+S}w\in N_c$. Write $d(t)=\sum_k d_k t^k$ for
its Laurent expansion at a primitive $p$th root in the constant
crystalline basis. Then there is a single integer $a_c$, independent
of $w$, and for each $w$ an $x\in N_{c-2}$ with
\[
 p^{a_c}d_0-\varphi_{\rm cris}^2\operatorname{cmp}(x)
   \in\operatorname{Fil}^{-\delta},\qquad
 a_c\le 2H/q^2+1+y-b/p+\lambda+3 .
\]
The comparison with de Rham periods at $c$ is understood.

Put $P=Q_1^yQ^\delta$. Since $Q^yN_{c-1}\subset
A_0\varphi N_{c-2}$ by the upper height bound at $c-1$, we
have $Pw\in R_{c-2}^{(2)}$. Write its coordinates in the basis
$\varphi^2 e_j$ and use the inverse integral two-step comparison
of T83A1. It gives integral $b_j(v)$ such that, at the root,
\[
 p^{s_2}P(t)d(t)=\sum_j b_j(v(t))\mu_j\pmod{t^{X_0}},
 \qquad \mu_j=\varphi_{\rm cris}^2\operatorname{cmp}(e_j).
\]
The translated basis itself is regular there because two Frobenius
substitutions land at zero. Now project to the quotient by
$\operatorname{Fil}^{-\delta}$, denoting projection by a bar.
For $k\ne0$, $(u^k-1)d_k\in\operatorname{Fil}^{-k}$ by the
filtered Taylor criterion for $(\gamma-1)w$. For all $k<b$ with
$k\ne0$, it follows that $\bar d_k=0$, using monotonicity for
$k\le\delta$ and the jump gap from $-b$ to $-\delta$ otherwise.
As $P$ has order $\delta$, the projected equation through order
$b+\delta-1$ therefore reads
\[
 p^{s_2}P\bar d_0\equiv\sum b_j\bar\mu_j\pmod{Q^{b+\delta}}.
\]
Here $Q$ has a simple zero at the branch and is interchangeable
with the Taylor parameter over the coefficient field. Moreover the
coordinate change $v=Q-p$ identifies $O[[v]]$ with $O[[Q]]$ by
$p$-adic completeness.

Choose an $O$-basis of the lattice generated by the $\bar\mu_j$.
The leading nonzero coefficient of $P$ at order $\delta$ shows
that $p^{s_2}\bar d_0$ belongs to its rational span. Each
coordinate $\theta$ of that vector satisfies
$\theta Q_1^y Q^\delta\equiv B(Q)\bmod Q^{b+\delta}$
for some $B\in O[[Q]]$. Cancelling $Q^\delta$ says the first
$b$ coefficients of $\theta Q_1^y$ are integral. Now
$Q_1=p\varepsilon+UQ^p$, where $\varepsilon=1-UQ^q$
and $U$ are units in $O[[Q]]$. The $j$th binomial summand of
$Q_1^y$ has initial degree $pj$ and coefficient valuation
$y-j+v_p\binom yj$, in the normalization $v_p(p)=1$.
If the entire truncation multiplied by $\theta$ is integral, each
summand with $pj<b$ must have integral scalar content: otherwise
take the smallest offending $j$ and inspect its first degree, since
all preceding summands have integral coefficients and later ones
start in larger degree. Take $j=\lfloor(b-1)/p\rfloor<y$,
using $1\le b\le py$ by thickness. The factorial valuation formula
bounds $v_p\binom yj\le\lfloor\log_p y\rfloor$ by counting
carries. Thus $-v_p(\theta)\le y-b/p+1+\log_p y$.
Taking $a_c=s_2+\lceil y-b/p+2+\log_p y\rceil$ works for
every coordinate. Its bound follows from $X_0\le2H$ and
$y\le H$. The basis vectors of the output lattice are integral
linear combinations of the $\bar\mu_j$; lifting the combination
to $N_{c-2}$ gives the asserted $x$. Tame conjugacy transports
the result to every primitive branch: it rescales $t$ while fixing
the constant coefficients. Multiplication by the clearing power of the nonzero series
$t$ times a unit is injective on lower-bounded Laurent series, so
any clearing exponent and expansion in the native test give this
same $d_0$. This proves the native constant-coefficient
assertion used in B8.42bue, Claim~1.
% END T83A2

% BEGIN T83A3
\paragraph{T83A3. The determinant loss after discarding the central generators.}
At a negative letter $i$ let $b=\beta_i+S$,
$R=A_0\varphi N_{i-1}$ and $X_i=vN_i+Q^bR\subset N_i$.
There are at most $m$ vectors $g_j\in N_i$ such that
\[
 \frac{\operatorname{length}_O(N_i/(X_i+\sum Og_j))}
 {\operatorname{length}_O(O/pO)}
 \le C_{\rm len}:=rS+2m\bigl((H+1)/q+1\bigr).
\]
Indeed choose a basis $e'$ of $N_{i-1}$ adapted, after specialization
and crystalline Frobenius, to the filtration steps from
$-\delta$ to $\delta+1$. Their intersections with $O^r$ are
saturated and hence have simultaneous adapted integral bases over
the DVR. The first $m'\le m$ vectors account for the central jumps
$u_j\in[-\delta,\delta]$, with multiplicity.
In coordinates at $i$ the determinant of $Q^b\varphi e'$ is
a unit times $Q^{w_0}$, where
$w_0=br+\sum_k k\,\operatorname{mult}_i(k)$.
The reverse height inclusion makes its determinant a divisor of
a power of $Q$. The Taylor lattice in W79H2--W79H4 gives the
order: in a filtration basis a period of jump $k$ contributes
$b+k$ to the determinant of these regular crystalline coordinates
relative to the Wach lattice.

Apply the first truncation of T83A1 with $(M,C)=(1,b+\delta)$.
The resulting columns $f_j\in R$ have constant specializations
$p^{s_i}\varphi\operatorname{cmp}(e'_j)$, where
$s_i=\lceil(1+b+\delta)/q+1\rceil$, and agree with these
constant periods through Taylor degree $b+\delta-1$ at the branch.
For $j\le m'$ put $g_j=Q^{-u_j}f_j$. This is integral:
the main term starts in degree $-u_j$ with values in
$\operatorname{Fil}^{u_j}$, and its error starts in degree
$b+\delta-u_j\ge b$, where every vector passes the filtration
test. All poles can first be cleared by $Q^{b+u_j}$, so
D82P1's integral divisibility proof applies.
Use these $g_j$ as the first columns of a matrix in $N_i$ and
$Q^b\varphi e'_j$ for the others. At $v=0$ its determinant has
normalized valuation
\[
       w_0+s_i m'-\sum_{j\le m'}(b+u_j)\le rS+s_i m' .
\]
Indeed multiplication of the first columns by $Q^{b+u_j}$ changes
the original matrix by its $p^{s_i}$ truncation on those columns
at zero. The uncancelled jumps belong to $[-b,-\beta_i]$ and
contribute at most $S$ apiece. Since $s_i\le2((H+1)/q+1)$,
the determinant bound on the cokernel in $N_i/vN_i$ gives
the claimed normalized length.
% END T83A3

% BEGIN T83A4
\paragraph{T83A4. The comparison on the neighbouring positive letter.}
Let $i$ be positive, $B=\beta_i$, $b'=\beta_{i-1}+S$,
$s_2=\lceil2H/q^2+1\rceil$ and
\[
 E_i=\left\lfloor(3a+S+2\delta-1)/q+\log_p H\right\rfloor .
\]
For $y\in N_{i-2}$ with $d=\varphi^2\operatorname{cmp}(y)
\in\operatorname{Fil}^B$, there is an integral invariant
source vector $z$ at $i$ such that
$\operatorname{cmp}(Fz)=p^{E_i+s_2+b'-B}d$.
The integer $s_2+b'-B$ need not be positive.

Use the two-step frame of T83A1 to represent the specialized
coordinates of $y$ by $l\in R_{i-2}^{(2)}$. Its value at zero is
$p^{s_2}d$; its Taylor series at each primitive root is the same
constant modulo $t^{2H}$. Thus
$x=Q^{-B}Q_1^{b'}l$ has Taylor terms of degrees below $2H-B$
all proportional to $d$ and none below $-B$. The ones of higher
degree satisfy the unrestricted part of the Taylor test, because
$2H-B\ge\delta$. Multiplication by $Q_1^{b'}$ puts $l$ in
$R_{i-1}^{(1)}$ by the lower height at $i-1$, and $Q^\delta$
clears the remaining height there. Hence a power of $Q$ clears
$x$, so its filtered Taylor test gives $x\in N_i$. At zero its
comparison is $p^{s_2+b'-B}d$.

Put $D=(\gamma-1)/v$, $c(j)=(1-u^{-j})/p$ and
$P_U(X)=\prod_{j\in U}(X-c(j))$. The integral operator $D$
preserves $N_i$ and $Q^{-a}N_i$: the ratio $\gamma(Q)/Q$ is
$1$ modulo $v$, as in L81S3. At a primitive branch, $v(0)=-p$
in the Taylor coordinate, so the diagonal entry on $t^k$ is
$c(-k)$. Take
\[
 I=[B-a+1,B],\qquad
 J=[B+1,B+S+a]\ \cup\ [1-\delta,\delta+a]
\]
with integral endpoints. For the element $x$ just constructed,
$P_I(D)x\in Q^aN_i$. Indeed on the $d$-line the factors successively
eliminate degrees $-B,\ldots,-B+a-1$. After dividing by $Q^a$ the
line still starts in degree $-B$. The error after division starts
in degree $2H-B-a\ge\delta$, because $a\le B\le h$.
The Taylor criterion again proves integral membership. By T83A1
both $P_I(D)x$ and $Q^a x$ lift to integral vectors in the
source. For a lift $y_0$ of $Q^a x$, $P_J(D)Q^{-a}y_0$ is
integral in the source: on each filtration basis line of jump $k$,
the pole-jet quotient $Q^{-a}N_i/N_i$ has diagonal values
$c(k+1),\ldots,c(k+a)$. Over the coefficient field these are
distinct on the line. The operator preserves every such line in the
Taylor model, and $J$ contains all its labels. Thus its product
kills the quotient there, and the integral Taylor criterion descends
membership. This gives a lift of $P_J(D)x$.

The two sets are disjoint by thickness. For $j\in I$ the valuation
of $P_J(c(j))P'_I(c(j))$ is $\sum_{k\in I\cup J,\,k\ne j}
v_p(j-k)$. The long interval containing $j$ contributes at most
$(2a+S-1)/q$ by the two factorial bounds. The other interval
contributes at most $(a+2\delta)/q+\log_p H$ by counting
multiples of each $p^e$; every nonzero difference within this
second count has absolute value at most $H$. Thus Lagrange
interpolation lifts $p^{E_i}$ to the polynomial ideal $(P_I,P_J)$
over $O$: interpolate the values $p^{E_i}/P_J(c(j))$ at the
roots of $P_I$, whose above derivative denominators are covered,
and divide the difference by the monic $P_I$. Applying these
integral polynomials to the two lifts gives a lift of $p^{E_i}x$,
as required.
% END T83A4

% BEGIN T83A5
\paragraph{T83A5. The three restrictions of the coefficient module.}
Let $I_F=FZ_f(T_1)+B(T_2)$. Lift the at most $(n+1)r$ generators
of $H_f^1(K,T_2)$ of I82B1 to cocycles. This defines an
$O$-linear map from a free module $P$ surjecting modulo boundaries.
A generator step below has quotient generated by at most $m$
elements, and a length step has normalized quotient length at most
$C_{\rm len}$. Put
\[
 M=\lfloor h/q\rfloor+1,\quad
 \sigma=(M+\delta)/q+1,\quad s=\lceil\sigma\rceil,\quad
 \sigma_2=2H/q^2+1,\quad
 E_G=(3a+S+2\delta)/q+\lambda.
\]
We apply the three constructions B8.42bud, B8.42bue and B8.42buf
of the companion on this $P$. We give their coefficient maps and
verify how their local inputs are supplied on the tested lattices.

\emph{Reduction of the cocycle.}
Write $D=(\gamma-1)/v$. Choose integral Herr representatives
$(y,z)$ on a basis of $P$. The crystalline extensions by $T_2$
and the $p^N,Q^a$ lifts of the trivial quotient from T83A1 give
differences $\mu_i,\nu_i$ between these two Wach lifts and
$p^N,Q^a$ times the same invariant \'etale lift. Thus
$p^N\nu_i-Q^a\mu_i\in N_i$,
$p^Nz_i+(\gamma-1)\mu_i\in vN_i$ and
$r_i=p^Ny_i+\varphi\mu_{i-1}-\mu_i$ lies in the localized
height intersection: $Q^{M_i}r_i\in N_i$, and
$p^{j_i}Q^{h_i}r_i\in R_{i-1}^{(1)}$ for some $j_i$,
where $[-M_i,h_i]$ are the local height intervals. Indeed the
first inclusion follows by the lower height on the extension.
For the second, lift the two Wach terms in the Frobenius span after
the upper height on the extension, and intersect with the kernel
after inverting $p$; A81R10 gives the kernel intersection there.
There are finitely many choices on a basis of $P$.

For a negative $c$ put $S_c=\{k+j\ne0:k\text{ a jump at }c,1\le j\le a\}$. A polynomial $R(D)$ with roots $c(t)$,
$t\in S_c$, and constant $p^{u_1}$ exists integrally for a
common $u_1\le(2a+2\delta+S)/q+\lambda$.
Indeed the central interval $[1-\delta,a+\delta]$ containing
zero contributes at most $(a+2\delta)/q$ by the factorial bound.
The remote interval $[1-\beta_c-S,a-\beta_c]$ contributes at most
$(a+S)/q+\log_p(2H)+1$ by counting multiples. By T83A4's
primitive pole-jet calculation $D R(D)$ takes $Q^{-a}N_c$
into $N_c$, including on the filtration lines with eigenvalue
zero. With $R_1=(R-R(0))/X$, take
\[
 m_c=R(D)Q^{-a}\nu_c+R_1(D)(z_c/v),\qquad
 u_c=p^Nm_c-p^{u_1}\mu_c,\qquad
 \zeta_c=p^{u_1}z_c+(\gamma-1)m_c.
\]
Then $Q^a u_c\in N_c$, $Du_c\in N_c$, and $\zeta_c\in vN_c$.
In fact $u_c=R(D)(p^NQ^{-a}\nu_c-\mu_c)
 +R_1(D)(p^Nz_c/v+D\mu_c)$ and
$p^N\zeta_c/v=p^{u_1}(p^Nz_c/v+D\mu_c)+Du_c$.
Cancel $p^N$ by saturation. Exactly as in T83A2, all the
nonconstant Laurent coefficients of $u_c$ satisfy the filtration
test because its gamma difference does. Its constant is in
$\operatorname{Fil}^{-a}$ since $Q^a u_c$ is integral.
Impose that it is in $\operatorname{Fil}^0$; this is one
generator step since the gap excludes all noncentral jumps in
$[-a,-1]$. Thus $u_c\in N_c$ on the new module.

For $z'\in vN_i$, solve $(\gamma-1)b=z'$ as a formal series
with zero constant in the invariant basis. The coefficient of
degree $k$ divides by $u^{qk}-1$ and consequently has denominator
bounded by $p^{w(k)}$, as in T83A1. Truncating $p^s\varphi b$
modulo $v^MQ^\delta R_i^{(1)}$ is integral by that claim.
Call it $t(z')$. This operation is linear and natural modulo that
submodule: the zero-constant solution is unique, and terms of
order $\ge M$ disappear since $\delta\le qM$. It vanishes there
for $z'\in v^MN_i$. In particular, writing
$t_c=t(\zeta_c)$ and
$\xi_c=p^s\varphi\zeta_c-(\gamma-1)t_c$, we have
$\xi_c\in v^MN_{c+1}$ by the lower height at the positive letter.
Set
\[
 \widehat m_c=p^s\varphi m_c+p^{u_1+s}y_{c+1}-t_c,\qquad
 \widehat u_c=p^N\widehat m_c-p^{u_1+s}\mu_{c+1}.
\]
The identity $\widehat u_c=p^s\varphi u_c+
p^{u_1+s}r_{c+1}-p^Nt_c$ gives
$Q^\delta\widehat u_c\in N_{c+1}$, and the Herr identity gives
$(\gamma-1)\widehat u_c\in N_{c+1}$. The constant-coefficient
kernel in $\operatorname{Fil}^{-\delta}/\operatorname{Fil}^0$
costs one further generator step and makes $\widehat u_c$ integral.

On these kernels put
$w_c=\varphi\widehat m_{c-2}+p^{u_1+s}y_c-p^sm_c$.
Its multiple $p^Nw_c=\varphi\widehat u_{c-2}-p^su_c+
p^{u_1+s}r_c$ is taken into $N_c$ by $Q^{\beta_c+S}$ and into
$R_{c-1}^{(1)}$ by $Q^\delta$ after a further scalar power of
$p$. Both lattices are $p$-saturated in the \'etale module (in
the translated basis for the second). Furthermore
$(\gamma-1)w_c=\varphi\xi_{c-2}-p^s\zeta_c\in vN_c$.
Thus $w_c$ satisfies T83A2. For any such almost invariant $w$,
or source vector with just finite cyclotomic poles and the same
gamma condition, use
$\tau(w)=p^s\varphi w-t((\gamma-1)w)$ as the first Herr entry.
Its gamma difference is in $v^MN_{c+1}$, and H82A1--H82A2
give a finite class by completing with $(\varphi-1)^{-1}$ on
$v^MN$. This inverse exists since $qM>h$. If $w\in vN$ then
$\tau(w)$ itself is in $v^MN$, since the solution with zero
constant for $(\gamma-1)w$ is $w$. Thus a congruence for $w$
modulo $vN$ gives the same congruence of classes, naturally for
$F$. For the vectors above, $\tau(w_c)\equiv
p^s(\varphi w_c+t_c)\bmod v^MN$. Directly,
\[
 \varphi w_c+t_c
 =p^{u_1+s}(y_{c+1}+\varphi y_c)
      +\varphi^2\widehat m_{c-2}-\widehat m_c.
\]
Summing over negative $c$ identifies the sum of first entries
with $p^{u_1+2s}y$ modulo a Frobenius boundary and $v^MN$.
The matching second entries follow by the Herr relation:
the boundary for $p^{u_1+s}(y,z)$ uses
$\sum_c(p^{u_1+s}y_c+\widehat m_c+\varphi\widehat m_c)$,
and its corrected second component is
$\sum_c(\xi_c+\varphi\xi_c)\in v^MN$. Injectivity of
$\varphi-1$ there finishes the class identity.
The total cost to this point is $n$ generator steps.

\emph{Matching the possibly nonintegral vector.}
Write $B=\beta_c+S$, $b_0=\beta_{c-2}+S$, and
$d_0=[t^0]w_c$. T83A2 gives $x_f\in N_{c-2}$ representing
$p^{a_c}d_0$ modulo $\operatorname{Fil}^{-\delta}$ through
$\varphi^2\operatorname{cmp}$. Its class modulo the inverse
image of that filtration is unique and linear in the coefficient
$f$. T83A3 at $c-2$ allows one length step and one generator
step to restrict the image to $X_{c-2}$ modulo the same ambiguity:
first require membership after adjoining the $Og_j$, then remove
these generators. Preimages preserve the two bounds, since
submodules of an $m$-generated module over a DVR are
$m$-generated. Specializing an element of $X_{c-2}$ gives
$p^{b_0}\varphi\operatorname{cmp}(t_f)$,
$t_f\in N_{c-3}$. Require $d_f=\varphi^2\operatorname{cmp}(t_f)$
to belong to $\operatorname{Fil}^{\beta_{c-1}}$ modulo the preceding
ambiguity; the possible obstruction lies in a quotient of the
projection of the specialization lattice to
$D_{c-1}/\operatorname{Fil}^{\beta_{c-1}}$, a vector space
of dimension at most $m$. That projected lattice is finite free
of rank at most $m$ (specialization uses $r$ constant coordinates).
Its quotient and the image submodule therefore need at most
$m$ generators. This is a second generator step, and choices on a basis of the
restricted coefficient module give a linear lift.

At $i=c-1$, T83A4 lifts $d_f$ at exponent
$E_i+\lceil\sigma_2\rceil+b_0-\beta_i$.
Apply the one-step comparison with $(M,C)=(1,B)$ to this source
lift at its constant specialization. This gives a vector of
finite poles with gamma difference in $vN_c$: indeed its defect
is in $vQ^B R_{c-1}^{(1)}\subset vN_c$. Its target Taylor
coefficients in degrees $<B$ equal a constant
$p^{k_c}\varphi d_f$, where
$k_c=E_i+\lceil\sigma_2\rceil+b_0-\beta_i+
\lceil(1+B)/q+1\rceil$. Combining with T83A2 uses only
$e_{\rm vec}\ge a_c+k_c-b_0$. The term $-\beta_c/p$
then combines with $B/q$ to leave $\beta_c/(pq)+S/q$.
Consequently one may take the common integer
\[
 e_{\rm vec}=\left\lfloor h/(pq)+(S+1)/q+S+
              E_G+2\sigma_2+6+\lambda\right\rfloor .
\]
Here $E_i\le E_G$, $a_c\le\sigma_2+\beta_{c-1}+S-
\beta_c/p+\lambda+3$, and each rounding error is less than
one. Scale the choice of $x_f$ by the excess integer in this
inequality. The constant difference between $p^{e_{\rm vec}}w_c$
and the resulting source image lies in $\operatorname{Fil}^{-\delta}$.
Require its image modulo $\operatorname{Fil}^0$ to vanish, a third
generator step. All nonzero coefficients below $B$ meet the
filtration condition by the gamma difference of $w_c$ and
vanishing of the source terms there. Higher degrees do so because
$-B$ is the lower jump. D82P1 gives an integral remainder
$r(f)\in N_c$. The source vector used in this correction has
finite poles and gamma difference in $vN_c$.

\emph{Matching the integral remainder.}
Apply T83A3 once at $c$ to write
$\operatorname{cmp}(r)=p^{\beta_c+S}\varphi d_1$
with $d_1\in\operatorname{cmp}(N_{c-1})$, after one length
and one generator step. The upper height at $c-1$ then gives
$p^{\beta_{c-1}+S}d_1=\varphi d_2$ with
$d_2\in\operatorname{cmp}(N_{c-2})$. Apply T83A3 a second
time, using a representative of $d_2$ modulo $vN_{c-2}$:
one further length and generator step give
$d_2=p^{b_0}\varphi d_3$, $d_3\in\operatorname{cmp}(N_{c-3})$.
Specialization kills exactly $vN$ on invariant components and all
three resulting period maps are linear by injectivity of Frobenius.
A third generator step imposes
$\varphi^2d_3\in\operatorname{Fil}^{\beta_{c-1}}$ as above.
T83A4 supplies a source lift $z$ with
$\operatorname{cmp}(Fz)=p^{E_i+\lceil\sigma_2\rceil+b_0-\beta_i}
 \varphi^2d_3$ for $i=c-1$.
Then $p^{\lceil E_G+\sigma_2\rceil-E_i-\lceil\sigma_2\rceil}
Q^{\beta_c+S}\varphi z$ lies in the source Wach lattice by the
lower height, and its specialization agrees after $F$ with that
of $p^{e_{\rm int}}r$, $e_{\rm int}=S+\lceil E_G+\sigma_2\rceil$.
Their difference lies in $vN_c$.
Functoriality of the Herr class just proved now lifts each
negative component. With $e=u_1+2s+e_{\rm vec}+e_{\rm int}$,
the image of the terminal coefficient module after $p^e$ lies
in $I_F$. On the whole word there were at most
$n+(3+3)n/2=4n$ generator steps and $(1+2)n/2=3n/2$ length steps.
% END T83A5

% BEGIN T83A6
\paragraph{T83A6. The actual frozen C and D row at every all-thick place.}
For each test in T83A1 the normalized length of
$Z_f(T_2)/I_F$ is at most the actual hub's $\mathrm{lenRowCD}(w)$,
for either band assignment. In particular the global middle
inequality holds at every guarded inert index with
$\psi=\mathrm{cc}$, including when the aggregate test of C82S2
does not apply.

First the witness of T83A5 satisfies the pinned numerical bound
\[
            e\le 32H/p^2+4S+10\delta/q+25\lambda .
\]
Indeed before substituting the bound on $a$, the sum of its four
terms is at most
\[
 (8a+8\delta+4S)/q+2h/q^2+6H/q^2+
 h/(pq)+2S+14+3/q+4\lambda .
\]
Substitute $a\le H(2p-1)/q^2+4+\lambda$ and $h\le H$.
The coefficient of $H$, multiplied by $p^2$, is bounded by
$8(p/q)^2(2+1/q)+8(p/q)^2+p/q<32$ since $q\ge10$.
The remaining constants are at most
$14+35/q+(4+8/q)\lambda\le25\lambda$, and $2+4/q\le4$.

Let $E$ be the coarse exponent on the \emph{same} lattice map
of N82I2, $E\le25h_B/p+3$. Map the chain of T83A5 into
$Z_f(T_2)/I_F$. Its initial term is the whole quotient. Each
generator step costs at most $m E$ in normalized length because
$p^E$ kills the whole quotient. Each length step costs at most
$C_{\rm len}$, and the terminal image is generated by at most
$(n+1)r$ elements (its preimage is a submodule of $P$) and is
killed by $p^e$. The required length is therefore finite and
bounded by
\[
       (n+1)r e+4nm E+(3/2)n\{rS+2m((H+1)/q+1)\}.
\]
This use of the proved $E$ suffices for the concrete hub row,
independently of the canonical universal symbol $e_{PC}$.

For clarity write $F_w$ for the original flat term of N82I2 and
$U=1+h_B/p+\log_p h_B$. At the given point,
$H\le h_B+\eta_3$, $\delta\le\eta_3\le h_B$,
$S\le S_c$, $1\le h_B\le2\exp(\mathrm{large}_2)$ and
$\mathrm{large}_2\ge1$. Thus
$\log_p(2H)\le7+\mathrm{large}_2$, since
$2H\le8\exp(\mathrm{large}_2)$ and $\log p\ge1$.
It follows from $F_w\ge(w_P+\eta_3)/p+S_c+\mathrm{large}_2$
that $e\le32h_B/p^2+225 F_w$. Moreover
$E\le25U$ and $(H+1)/q+1\le4U$, using $p\ge11$.
Hence the three contributions are bounded in sum by
$32\,\mathrm{coreH}+227\,\mathrm{coreF}
 +112\,\mathrm{coreM}$. The frozen coefficients
$(368,2082,22365)$ dominate these numbers. The lower side of
the hub's rational rounding gives exactly the asserted per-map
length, with denominator $\operatorname{length}_O(O/pO)$.
The argument uses the actual coefficient ring throughout.

If $\psi=\mathrm{cc}$, C82S1 leaves only all-thick or all-nonthick
places. Apply the bound just proved and N82I1 respectively at
each place, then G82I1--G82I2 at the common point. This removes
the extra sufficient inequality in C82S2. For $\psi\ne\mathrm{cc}$
the remaining words may have both kinds of letter; their needed
bound is the C or D colength on those tested maps. The argument
here establishes the actual all-thick local test and its indicated
global specialization.
% END T83A6
